%==============================================================================
% CAPÍTULO 23 — RAG ODONTOLÓGICO COM LITERATURA CIENTÍFICA
% PARTE IV: APLICAÇÕES CLÍNICAS DA IA
%==============================================================================
\chapter{RAG Odontológico com Literatura Científica}
\label{cap:rag-odontologico}
\nivelquatro

\begin{quote}
\textit{``Um LLM sem acesso à literatura científica é como um dentista sem acesso a artigos: pode ter boa formação básica, mas suas recomendações estarão inevitavelmente desatualizadas e descontextualizadas. O RAG --- Retrieval-Augmented Generation --- resolve isso ancorando cada resposta em evidências recuperadas em tempo real.''}
\end{quote}

Um grande modelo de linguagem sem acesso à literatura científica contemporânea é como um dentista sem acesso a artigos: tem boa formação, mas seu conhecimento inevitavelmente se desatualiza e suas recomendações podem ser imprecisas. O RAG (Retrieval-Augmented Generation) resolve este problema ao ancorar cada resposta em evidências recuperadas em tempo real de uma base de conhecimento externa --- artigos do PubMed, diretrizes clínicas, revisões sistemáticas. Este capítulo apresenta a arquitetura completa do RAG, desde a indexação de documentos e geração de embeddings até a recuperação semântica e geração condicionada, com foco em aplicações odontológicas. A citação que abre este capítulo captura com precisão o espírito desta tecnologia: a síntese entre o conhecimento generalista dos LLMs e a precisão factual da literatura científica, tema que exploraremos em profundidade nas seções seguintes.

\subsection{Objetivos de Aprendizagem}

Ao final deste capítulo, você será capaz de:

\subsubsection{Compreender o Problema do Conhecimento Estático em LLMs}
Você entenderá por que modelos de linguagem puros são insuficientes para responder perguntas clínicas atualizadas sem acesso a bases externas.

\subsubsection{Dominar a Arquitetura RAG: Indexação, Recuperação e Geração}
Exploraremos cada componente do RAG --- embedding, banco vetorial, busca semântica e prompt aumentado --- com foco odontológico.

\subsubsection{Construir uma Base Vetorial com Artigos Odontológicos}
Você implementará um pipeline de ingestão de PDFs científicos com extração de trechos, metadados e DOIs para rastreabilidade.

\subsubsection{Controlar Alucinações e Garantir Evidência nas Respostas}
Aprenderá técnicas para que o sistema só responda com base em fontes recuperadas, reduzindo riscos de informação incorreta.

\section{Introdução: O Problema do Conhecimento Estático}

Os \textit{Large Language Models} (LLMs) possuem uma limitação fundamental: seu conhecimento é estático, congelado no momento do treinamento. O GPT-4, lançado em 2023, desconhece qualquer descoberta odontológica publicada após sua data de corte. Um clínico que perguntasse a um LLM não-aumentado sobre a nova classificação de peri-implantite de 2025 receberia informações de 2023 ou anterior --- potencialmente obsoletas ou incorretas.

Além da obsolescência, há o problema da alucinação (discutido no Capítulo 22): LLMs podem gerar informações plausíveis mas factualmente incorretas, especialmente em domínios especializados como a odontologia, onde o conhecimento é altamente técnico e a literatura é volumosa ($> 30.000$ artigos odontológicos publicados anualmente no PubMed).

O \textbf{Retrieval-Augmented Generation} (RAG), introduzido por Lewis et al. (Meta AI, 2020), resolve ambos os problemas com uma arquitetura elegante: antes de gerar uma resposta, o sistema recupera documentos relevantes de uma base de conhecimento externa e condiciona a geração a estas evidências.

\begin{figure}[H]
\centering
\begin{tikzpicture}[node distance=2.0cm, auto]
    \node [draw, fill=corPrimaria!10, rounded corners, text width=3cm, align=center] (query) {Pergunta do\\Usuário};
    \node [draw, fill=corSecundaria!10, rounded corners, text width=3.5cm, align=center, right=of query] (retrieval) {Busca Semântica\\\footnotesize{Embeddings + Vetorial DB}};
    \node [draw, fill=corDestaque!10, rounded corners, text width=3.5cm, align=center, right=of retrieval] (docs) {Documentos\\Recuperados\\\footnotesize{Artigos, guidelines}};
    \node [draw, fill=corNivelPhD!10, rounded corners, text width=4cm, align=center, below=of docs] (llm) {LLM + Prompt\\\footnotesize{Pergunta + Evidências $\rightarrow$ Resposta}};
    \node [draw, fill=corPrimaria!10, rounded corners, text width=3cm, align=center, below=of llm] (resp) {Resposta\\Fundamentada};
    \draw [->, thick] (query) -- (retrieval);
    \draw [->, thick] (retrieval) -- (docs);
    \draw [->, thick] (docs) -- (llm);
    \draw [->, thick] (llm) -- (resp);
    \draw [->, thick, dashed] (query) to [bend right] (llm);
\end{tikzpicture}
\caption{Arquitetura conceitual do RAG: a pergunta do usuário dispara uma busca semântica em base vetorial. Os documentos recuperados são injetados no prompt do LLM, que gera uma resposta ancorada nas evidências.}
\label{fig:rag-arquitetura}
\end{figure}


% === FIGURAS ADICIONADAS ===
% ======================================================================
% FIGURA: Arquitetura RAG Odontológico (Capítulo 23)
% ======================================================================
\begin{figure}[H]
\centering
\begin{tikzpicture}[
    node distance=1.5cm and 2cm,
    box/.style={rectangle, draw=blueaccent, fill=blueaccent!15, align=center, minimum width=2.5cm, align=center, minimum height=0.8cm, font=\footnotesize, rounded corners=3pt, align=center},
    db/.style={cylinder, draw=cyanaccent, fill=cyanaccent!15, align=center, minimum width=2cm, align=center, minimum height=1.2cm, font=\footnotesize, shape border rotate=90, align=center},
    arrow/.style={-latex, thick},
]
% Usuário
\node[align=center, box, fill=codegreen!20] (user) {Usuário\\Pergunta clínica};

% Retrieval
\node[align=center, box, right=of user] (embed) {Embedding\\Transformer\\768-d};
\node[align=center, db, right=of embed] (vectordb) {Vector DB\\100 artigos\\ChromaDB};

% Generation
\node[align=center, box, below=of embed] (retrieve) {Top-K\\Recuperação\\k=5};
\node[align=center, box, right=of retrieve] (llm) {LLM\\GPT-4/LLaMA\\Geração};
\node[align=center, box, right=of llm, fill=codegreen!20] (answer) {Resposta\\com DOI\\+ Confiança};

% Setas
\draw[arrow] (user) -- (embed) node[midway, above, font=\tiny] {Query};
\draw[arrow] (embed) -- (vectordb) node[midway, above, font=\tiny] {similaridade};
\draw[arrow] (vectordb) -- (retrieve) node[midway, right, font=\tiny] {trechos};
\draw[arrow] (retrieve) -- (llm) node[midway, above, font=\tiny] {contexto};
\draw[arrow] (llm) -- (answer) node[midway, above, font=\tiny] {resposta};
\draw[arrow, bend left=20] (answer.south) to (user.south);

% Título
\node[font=\small\bfseries, above=0.5cm of embed] {RAG Odontológico: Recuperação + Geração};

\end{tikzpicture}
\caption{Arquitetura de Retrieval-Augmented Generation (RAG) para consulta odontológica. A pergunta do usuário é convertida em embedding (768-d), comparada com a base vetorial de 100 artigos científicos, e os 5 trechos mais relevantes são usados como contexto pelo LLM para gerar uma resposta com citação DOI.}
\label{fig:rag-architecture}
\end{figure}


\section{Arquitetura do RAG: Componentes Fundamentais}

\subsection{Indexação: Construindo a Base de Conhecimento Vetorial}

O primeiro passo do RAG é transformar uma coleção de documentos (artigos científicos, diretrizes clínicas, livros-texto) em uma base de dados vetorial pesquisável semanticamente. O processo de indexação envolve:

\begin{enumerate}
    \item \textbf{Coleta de documentos}: Download de artigos em PDF ou extração de metadados (título, resumo, corpo do texto) via APIs como PubMed Entrez, Semantic Scholar ou Scopus. Idealmente, inclui-se o texto completo, não apenas o resumo;

    \item \textbf{Segmentação (\textit{chunking})}: Os documentos são divididos em segmentos (\textit{chunks}) de tamanho adequado --- tipicamente 512--1024 \textit{tokens} (aproximadamente 350--700 palavras). Este tamanho equilibra contexto suficiente para compreensão com especificidade para recuperação precisa. Estratégias de segmentação incluem:
    \begin{itemize}
        \item \textbf{Fixa}: Divisão a cada $N$ tokens, independente da estrutura do texto;
        \item \textbf{Estrutural}: Divisão por seções do artigo (Introdução, Métodos, Resultados, Discussão);
    \item \textbf{Deslizante (\textit{sliding window})}: Sobreposição de 10--20\% entre chunks adjacentes para evitar cortar informações no limite.
\end{itemize}

\subsubsection{Geração de Embeddings e Armazenamento}\indice{embeddings}\indice{banco vetorial}\indice{ChromaDB}

    \item \textbf{Geração de embeddings}: Cada chunk é convertido em um vetor denso de dimensionalidade fixa (ex.: 768 para BERT-based, 1536 para Ada-002 da OpenAI) usando um modelo de \textit{embedding}. O embedding captura o significado semântico do texto, de modo que textos com significado similar tenham vetores próximos no espaço;

    \item \textbf{Armazenamento em banco vetorial}: Os embeddings são armazenados em um banco de dados vetorial (\textit{vector database}) --- como Pinecone, Weaviate, Qdrant, Milvus ou ChromaDB --- que suporta busca por similaridade de cosseno (ou distância euclidiana) em escala.
\end{enumerate}

\subsection{Recuperação: Busca Semântica com Embeddings}

Quando o usuário faz uma pergunta (ex.: ``Qual a evidência para regeneração óssea guiada com membrana de PTFE vs. colágeno em defeitos peri-implantares?''), o sistema:

\subsubsection{Processo de Recuperação}\indice{busca semântica}\indice{re-ranking}

\begin{enumerate}
    \item Converte a pergunta em um embedding usando o mesmo modelo da indexação;
    \item Busca no banco vetorial pelos $k$ chunks mais similares (tipicamente $k = 5$--$20$);
    \item Recupera os textos completos dos chunks e seus metadados (DOI, título do artigo, autores, ano, seção do artigo);
    \item Opcionalmente, aplica \textbf{re-ranking}: um modelo mais preciso (mas mais lento) reordena os chunks recuperados por relevância, refinando a seleção inicial.
\end{enumerate}

\destaque{
\textbf{Métricas de similaridade comuns em bancos vetoriais:}
\begin{itemize}
    \item \textbf{Similaridade de cosseno}: $\cos(\theta) = \frac{\mathbf{u} \cdot \mathbf{v}}{\|\mathbf{u}\| \|\mathbf{v}\|}$. Varia de $-1$ (opostos) a $+1$ (idênticos). É a métrica mais comum, pois normaliza a magnitude dos vetores;
    \item \textbf{Distância euclidiana}: $d = \|\mathbf{u} - \mathbf{v}\|$. Menor distância = mais similar. Sensível à magnitude dos vetores;
    \item \textbf{Produto escalar}: $\mathbf{u} \cdot \mathbf{v}$. Equivalente à similaridade de cosseno quando vetores são normalizados.
\end{itemize}
}

\subsection{Geração: Condicionando o LLM às Evidências}

A etapa final insere os chunks recuperados no \textit{prompt} do LLM. O formato típico é:

\begin{lstlisting}[language={}, caption={Template de prompt para RAG odontológico}]
Voce e um assistente odontologico baseado em evidencias.
Responda a pergunta do usuario APENAS com base nos trechos
de artigos cientificos fornecidos abaixo. Cite a fonte (DOI)
de cada informacao. Se os trechos nao contiverem informacao
suficiente, diga explicitamente: "Nao ha evidencias suficientes
nos artigos consultados para responder esta pergunta."

## Evidencias recuperadas:

[TRECHO 1] (Fonte: doi:10.xxxx/yyyy, Secao: Resultados)
...texto do artigo...

[TRECHO 2] (Fonte: doi:10.xxxx/zzzz, Secao: Discussao)
...texto do artigo...

[TRECHO 3] (Fonte: doi:10.xxxx/wwww, Secao: Conclusao)
...texto do artigo...

## Pergunta do usuario:
{pergunta}

## Resposta fundamentada (com citacoes de DOI):
\end{lstlisting}

\subsubsection{Salvaguardas contra Alucinação}\indice{prompt engineering}\indice{citação obrigatória}

Este design de prompt implementa três salvaguardas contra alucinação:
\begin{enumerate}
    \item \textbf{Restrição explícita}: ``Responda APENAS com base nos trechos fornecidos'';
    \item \textbf{Citação obrigatória}: ``Cite a fonte (DOI) de cada informação'';
    \item \textbf{Reconhecimento de ignorância}: ``Se os trechos não contiverem informação suficiente, diga explicitamente...'' --- forçando o modelo a admitir quando não sabe, em vez de inventar.
\end{enumerate}

\section{Construindo uma Base Vetorial para Artigos Odontológicos}

\subsection{Seleção de Fontes: O Que Indexar?}

A qualidade do RAG depende criticamente da qualidade e abrangência da base documental. Para um RAG odontológico de alto desempenho, recomenda-se indexar:

\begin{itemize}
    \item \textbf{Artigos de revisão sistemática e metanálises} (PubMed, Cochrane Oral Health): Representam o mais alto nível de evidência. Tipicamente 2.000--5.000 artigos relevantes;

    \item \textbf{Diretrizes clínicas} (ADA, AAPD, AAP, EFP, AAO): Documentos oficiais de sociedades de especialidade. Aproximadamente 200--500 diretrizes ativas;

    \item \textbf{Ensaios clínicos randomizados} (PubMed, ICTRP): Evidência de eficácia de intervenções. 10.000--20.000 ECRs odontológicos publicados;

    \item \textbf{Artigos originais de alta qualidade}: Selecionados por periódico (Qualis A1--A2, ou fator de impacto JCR $>3$). Foco nos últimos 10 anos, com marcos históricos incluídos seletivamente.
\end{itemize}

\subsection{Metadados Enriquecedores}

Cada chunk indexado deve carregar metadados que permitam filtragem e avaliação crítica:

\begin{lstlisting}[language=Python, caption={Estrutura de metadados para chunks odontológicos}]
metadata = {
    "doi": "10.1016/j.jdent.2024.104456",
    "title": "Caries detection on bitewing radiographs...",
    "authors": ["Cantu, J.", "Lee, P.", "Schwendicke, F."],
    "year": 2024,
    "journal": "Journal of Dentistry",
    "qualis": "A1",
    "study_type": "original",
    "section": "Results",           # Secao do artigo de onde veio o chunk
    "evidence_level": "II",          # Nivel de evidencia (I-V)
    "specialty": ["cariology", "radiology"],
    "chunk_index": 7,                # Posicao do chunk no documento
    "total_chunks": 24               # Total de chunks do documento
}
\end{lstlisting}

\section{Recuperação de Trechos com DOI: Rastreabilidade como Requisito}

\subsection{O Princípio da Evidência Rastreável}

Em odontologia baseada em evidências, a rastreabilidade é um requisito não-negociável. Toda afirmação clínica deve ser lastreada por uma fonte verificável. O RAG implementa este princípio arquiteturalmente: como cada chunk recuperado carrega seu DOI, a resposta gerada pode --- e deve --- citar precisamente de onde cada informação foi extraída.

\destaque{
\textbf{Formato de citação RAG:} ``Estudos recentes indicam que a U-Net supera métodos tradicionais de limiarização na detecção de perda óssea periodontal, com sensibilidade de 87,5\% \textsuperscript{[1]}. A inclusão de dados demográficos (tabagismo, diabetes) aumenta a AUC em 5\% comparado ao uso exclusivo de imagens \textsuperscript{[2]}.''

\noindent\textbf{Referências geradas automaticamente:}
\begin{enumerate}
    \item Lee et al. (2024). Deep learning for detection of periodontal bone loss. \textit{J Clin Periodontol}. DOI: 10.1111/jcpe.13890
    \item Costa et al. (2024). ML models for periodontitis prediction. \textit{Clin Oral Investig}. DOI: 10.1007/s00784-024-05678-3
\end{enumerate}
}

\subsection{Controle de Qualidade da Recuperação}

\subsubsection{Estratégias de Filtragem}\indice{filtragem}\indice{nível de evidência}\indice{deduplicação}

Nem todo chunk recuperado é relevante ou de alta qualidade. Estratégias de filtragem incluem:
\begin{itemize}
    \item \textbf{Limiar de similaridade}: Descartar chunks com similaridade de cosseno $< 0,70$ (baixa relevância semântica);
    \item \textbf{Filtro de nível de evidência}: Priorizar revisões sistemáticas (nível I) e ECRs (nível II) sobre séries de casos (nível IV) e opiniões de especialistas (nível V);
    \item \textbf{Filtro de atualidade}: Priorizar artigos dos últimos 5 anos, exceto quando marcos históricos são solicitados;
    \item \textbf{Deduplicação semântica}: Se múltiplos chunks do mesmo artigo são recuperados, manter apenas o de maior similaridade;
    \item \textbf{Diversidade de fontes}: Garantir que a resposta sintetize múltiplos artigos independentes, não dependa de uma única fonte.
\end{itemize}

\section{Controle de Alucinação: Evidência como Requisito Mínimo}

\subsection{Tipos de Alucinação em RAG Odontológico}

\subsubsection{Categorias de Alucinação em RAG}\indice{alucinação}\indice{extrapolação}\indice{citação}

Mesmo com RAG, alucinações podem ocorrer. É útil categorizá-las para desenvolver contramedidas específicas:

\begin{enumerate}
    \item \textbf{Alucinação de extrapolação}: O LLM extrapola além do que os chunks dizem. Ex.: chunk afirma ``U-Net detecta cárie com sensibilidade 89\% em bitewing'' e LLM responde ``U-Net detecta cárie com sensibilidade 89\% em \textbf{todas as modalidades radiográficas}'';

    \item \textbf{Alucinação de síntese incorreta}: O LLM combina informações de múltiplos chunks de forma logicamente incorreta. Ex.: chunk A diz ``Clorexidina 0,12\% reduz placa'' e chunk B diz ``Listerine reduz gengivite''; LLM conclui erroneamente ``Clorexidina 0,12\% + Listerine é superior a Clorexidina isolada'';

    \item \textbf{Alucinação de citação}: O LLM inventa um DOI ou atribui uma afirmação ao artigo errado. Esta é particularmente perigosa porque aparenta rastreabilidade;

    \item \textbf{Alucinação de omissão}: O LLM omite \textit{caveats} importantes presentes nos chunks. Ex.: chunk menciona que o estudo ``excluiu pacientes fumantes'', mas o LLM omite esta limitação ao reportar os resultados.
\end{enumerate}

\subsection{Estratégias de Mitigação Específicas para RAG}

\begin{enumerate}
    \item \textbf{Verificação de \textit{grounding}}: Um segundo LLM (ou o mesmo, com prompt diferente) verifica se cada afirmação na resposta está de fato presente em pelo menos um dos chunks recuperados. Afirmações não-ancoradas são removidas ou sinalizadas;

    \item \textbf{Verificação de DOI}: O DOI citado é verificado contra os metadados dos chunks recuperados. DOIs que não constam na lista de chunks são marcados como potencialmente alucinados;

    \item \textbf{\textit{Self-consistency}}: Gerar três respostas independentes (diferentes temperaturas ou diferentes LLMs) e verificar consistência fática. Discordâncias indicam incerteza ou alucinação;

    \item \textbf{Limiar de confiança}: Se o modelo de embedding indica que a similaridade máxima entre a pergunta e qualquer chunk é $< 0,60$, o sistema deve responder: ``Não encontrei evidências suficientes na base de conhecimento para responder esta pergunta com confiança'';

    \item \textbf{Preservação de \textit{caveats}}: O prompt inclui instrução explícita para preservar limitações metodológicas dos estudos citados (tamanho de amostra, critérios de exclusão, validade externa).
\end{enumerate}

\section{Limites Éticos de Recomendação Clínica por LLM}

\subsection{O Princípio da Não-Substituição}

Um RAG odontológico, por mais sofisticado que seja, \textbf{não pode e não deve substituir o julgamento clínico do dentista}. Este princípio deve ser operacionalizado em múltiplos níveis:

\begin{enumerate}
    \item \textbf{\textit{Disclaimer} obrigatório}: Toda resposta deve iniciar com um aviso explícito: ``\textbf{Aviso:} Este sistema é uma ferramenta de suporte à decisão baseada em literatura científica. As informações fornecidas não constituem recomendação clínica individualizada. O diagnóstico e plano de tratamento devem ser determinados pelo dentista responsável, considerando o quadro clínico completo do paciente.''

    \item \textbf{Recusa de perguntas de prescrição}: O sistema deve recusar-se a responder perguntas como ``Qual a dose de amoxicilina para profilaxia de endocardite?'' ou ``Quantos tubetes de lidocaína 2\% posso usar neste paciente?'' --- estas são decisões que dependem de avaliação individual (peso, alergias, função renal/hepática) e devem ser respondidas pelo dentista, não por um LLM;

    \item \textbf{Recusa de diagnóstico individual}: ``Este paciente tem periodontite estágio III grau B?'' --- o sistema pode descrever os critérios diagnósticos e citar a classificação, mas não pode aplicar os critérios a um caso individual sem acesso ao exame clínico completo e responsabilidade profissional;

    \item \textbf{Escalada para especialista}: Para perguntas sobre procedimentos de alta complexidade (cirurgia ortognática, implantes zigomáticos, enxertos microvasculares), o sistema deve recomendar consulta com especialista, não tentar fornecer orientação detalhada.
\end{enumerate}

\subsection{Responsabilidade e Arcabouço Regulatório}

No Brasil, a Resolução CFO-196/2019 (que dispõe sobre publicidade e uso de tecnologias na odontologia) e o Código de Ética Odontológica estabelecem que a responsabilidade pelo diagnóstico e tratamento é exclusiva do cirurgião-dentista. Sistemas de IA, incluindo RAG, são instrumentos auxiliares. A ANVISA e o Conselho Federal de Odontologia ainda não publicaram regulação específica para \textit{software} de suporte à decisão clínica baseado em IA, mas é questão de tempo.

\section{Prática: Construindo um Mini-RAG com Literatura Odontológica}

\begin{pratica}{Mini-RAG Odontológico com PubMed e ChromaDB}
Nesta prática, construímos um sistema RAG mínimo --- porém funcional --- que indexa artigos odontológicos do PubMed, armazena embeddings em ChromaDB e responde perguntas com evidências.

\begin{lstlisting}[language=Python, caption={Mini-RAG odontológico: indexação e consulta}]
from odontoia.rag import DentalRAG, PubMedLoader
from odontoia.rag.embedders import OpenAIEmbedder
from odontoia.rag.vectorstores import ChromaVectorStore

# 1. Carregar artigos do PubMed
loader = PubMedLoader(
    query="(dental OR oral) AND (deep learning OR AI) AND 2020:2025",
    max_results=200,
    include_fulltext=True   # Tenta baixar texto completo via PMC
)
documents = loader.load()
print(f"Carregados {len(documents)} artigos")
# -> Carregados 187 artigos (alguns sem fulltext disponivel)

# 2. Indexar em banco vetorial
embedder = OpenAIEmbedder(
    model="text-embedding-3-small",
    dimensions=1536
)
vectorstore = ChromaVectorStore(
    collection_name="odontologia_rag",
    persist_directory="./chroma_db",
    metric="cosine"
)

rag = DentalRAG(
    embedder=embedder,
    vectorstore=vectorstore,
    chunk_size=512,       # tokens por chunk
    chunk_overlap=64,     # sobreposicao entre chunks
    metadata_fields=[     # Metadados obrigatorios
        "doi", "title", "authors", "year",
        "journal", "study_type", "evidence_level"
    ]
)

# 3. Indexar (com barra de progresso)
rag.index(documents, show_progress=True)
print(f"Total de chunks indexados: {rag.vectorstore.count()}")

# 4. Consultar o RAG
pergunta = (
    "Qual a evidencia para uso de deep learning na deteccao "
    "de carie em radiografias panoramicas?"
)

resposta = rag.query(
    question=pergunta,
    k_retrieval=8,            # Recuperar 8 chunks
    k_final=5,                # Usar top-5 apos re-ranking
    min_similarity=0.65,      # Limiar minimo de similaridade
    evidence_level_filter=["I", "II"],  # Priorizar meta-analises e ECRs
    max_age_years=5,          # Apenas artigos de 2020-2025
    require_citations=True,   # Exigir citacoes com DOI
    disclaimer=True           # Incluir disclaimer etico
)

print("\n=== RESPOSTA ===")
print(resposta.answer)
print("\n=== EVIDENCIAS (ordenadas por relevancia) ===")
for i, ev in enumerate(resposta.evidence):
    print(f"\n[{i+1}] {ev.title}")
    print(f"    DOI: {ev.doi} | Ano: {ev.year}")
    print(f"    Similaridade: {ev.similarity:.2f}")
    print(f"    Trecho relevante: \"{ev.snippet[:200]}...\"")

# 5. Avaliar qualidade da resposta
print(f"\n=== METRICAS DE QUALIDADE ===")
print(f"Similaridade media: {resposta.mean_similarity:.2f}")
print(f"Fontes unicas: {len(set(e.doi for e in resposta.evidence))}")
print(f"Alucinacao provavel: {resposta.hallucination_risk}")
print(f"Nivel de evidencia predominante: {resposta.dominant_evidence_level}")
\end{lstlisting}

\textbf{Resultado esperado:}
\begin{lstlisting}[style=saida]
Carregados 187 artigos
Indexando... [####################] 100% (2,143 chunks)
Total de chunks indexados: 2,143

=== RESPOSTA ===
[Aviso: Este sistema e uma ferramenta de suporte a decisao...]

A literatura recente demonstra que modelos de deep learning,
particularmente U-Net e YOLO, alcancam alta acuracia na deteccao
de caries em radiografias panoramicas. Uma metanalise de 22
estudos [1] reportou sensibilidade agrupada de 89.5% e
especificidade de 91.3%. Modelos baseados em YOLO mostraram
melhor relacao velocidade-acuracia para aplicacoes em tempo
real [2]. Lesoes proximais sao detectadas com maior acuracia
que lesoes oclusais [1]. A principal limitacao e a qualidade
variavel dos datasets de treinamento [3].

=== EVIDENCIAS (ordenadas por relevancia) ===
[1] Lee et al. (2024) - Deep learning for caries detection in
    panoramic radiographs: a systematic review and meta-analysis
    DOI: 10.1177/0022034524123456 | Ano: 2024
    Similaridade: 0.92
    Trecho: "Pooled sensitivity of 89.5% (95%CI: 86.2-92.1%)
    and specificity of 91.3%..."

[2] Almeida et al. (2024) - Performance of deep learning models
    for caries detection in panoramic radiography
    DOI: 10.1259/dmfr.20230456 | Ano: 2024
    Similaridade: 0.88
    Trecho: "YOLO-based models demonstrated superior speed-
    accuracy trade-off for real-time applications..."

[3] Zhang et al. (2023) - Children's dental panoramic radiographs
    dataset for caries segmentation
    DOI: 10.1038/s41597-023-02319-6 | Ano: 2023
    Similaridade: 0.84
    Trecho: "The dataset provides a public benchmark for
    reproducible research..."

=== METRICAS DE QUALIDADE ===
Similaridade media: 0.88
Fontes unicas: 3
Alucinacao provavel: BAIXA (grounding score: 0.94)
Nivel de evidencia predominante: I (metanalises/revisoes)
\end{lstlisting}
\end{pratica}

\teste{O RAG deve indexar pelo menos 2.000 chunks a partir de 150+ artigos. Em consultas de teste, deve recuperar pelo menos 3 fontes distintas com similaridade $\ge 0,80$ e reportar grounding score $\ge 0,85$. O sistema deve recusar perguntas de prescrição e diagnóstico individual com mensagem apropriada.}

\section{Métricas de Avaliação de RAG}

Avaliar a qualidade de um sistema RAG é mais complexo que avaliar um classificador (onde basta acurácia/AUC). As dimensões de avaliação incluem:

\begin{enumerate}
    \item \textbf{Fidelidade (\textit{faithfulness} ou \textit{grounding})}: A resposta gerada é fiel aos chunks recuperados? Ou o LLM extrapolou/adicionou informações? Medida via \textit{grounding score} (0--1), tipicamente calculada por um LLM verificador;

    \item \textbf{Relevância da resposta (\textit{answer relevance})}: A resposta endereça diretamente a pergunta? Medida via similaridade de cosseno entre embedding da pergunta e embedding da resposta;

    \item \textbf{Relevância do contexto (\textit{context relevance})}: Os chunks recuperados são relevantes para a pergunta? Medida pela similaridade média entre pergunta e chunks;

    \item \textbf{Precisão de citação}: As citações (DOI) estão corretas? O trecho citado realmente contém a informação atribuída? Medida via verificação automática ou manual;

    \item \textbf{Cobertura}: A base de conhecimento contém informações para responder a gama esperada de perguntas? Medida via \textit{benchmark} de perguntas de referência com respostas conhecidas.
\end{enumerate}

\begin{table}[H]
\centering
\caption{Métricas de Avaliação de RAG Odontológico}
\label{tab:rag-metrics}
\begin{tabular}{p{3cm}p{2cm}p{2.5cm}p{4cm}}
\toprule
\textbf{Métrica} & \textbf{Alvo} & \textbf{Método} & \textbf{Interpretação} \\
\midrule
Grounding score & $\ge 0,85$ & LLM verificador & Proporção de afirmações ancoradas nos chunks \\
Answer relevance & $\ge 0,80$ & Similaridade cosseno & Resposta aborda a pergunta? \\
Context relevance & $\ge 0,75$ & Similaridade média & Chunks são pertinentes? \\
Citation accuracy & 100\% & Verificação DOI & Citações são corretas? \\
Coverage & $\ge 80\%$ & Benchmark & Base cobre as perguntas esperadas? \\
\bottomrule
\end{tabular}
\end{table}

\section{Estudos de Caso: RAG em Cenários Odontológicos Reais}

\subsection{Caso 1: Suporte à Decisão em Periodontia com RAG PubMed}

Em uma clínica de periodontia universitária, 15 residentes utilizaram um sistema RAG odontológico (indexando 8.500 artigos periodontais do PubMed, 2018--2025) como ferramenta de suporte durante o planejamento de casos complexos. Durante 6 meses e 320 consultas, o RAG foi consultado em 68\% das sessões de planejamento.

\subsubsection{Cenários de Uso e Avaliação}\indice{periodontia}\indice{RAG}\indice{suporte à decisão}

Os cenários de uso mais frequentes foram:
\begin{itemize}
    \item \textbf{Seleção de biomaterial para regeneração}: ``Qual a evidência para uso de rhPDGF-BB vs. PRF na regeneração de defeitos infraósseos de 3 paredes?'' --- o RAG recuperou 5 artigos (3 ECRs, 2 revisões sistemáticas) sintetizando as evidências comparativas;

    \item \textbf{Antibioticoterapia profilática}: ``Há evidência para antibioticoterapia profilática em raspagem subgengival de pacientes diabéticos controlados (HbA1c $<7\%$)?'' --- o RAG concluiu que a evidência é insuficiente (apenas 2 estudos, ambos com $n < 50$), e recomendou decisão individualizada;

    \item \textbf{Intervalo de manutenção}: ``Qual o intervalo ótimo de manutenção periodontal para pacientes com periodontite estágio III grau B?'' --- o RAG recuperou diretrizes da EFP (2021) e 3 ECRs, recomendando intervalo de 3--4 meses.
\end{itemize}

A avaliação dos residentes sobre o sistema (escala Likert 1--5) foi positiva: utilidade clínica (4,3/5), confiança nas respostas (4,1/5), facilidade de uso (4,5/5). O \textit{grounding score} médio das respostas foi de 0,91 (DP 0,06), indicando alta fidelidade aos artigos recuperados. Em apenas 3 ocasiões ($<1,5\%$ das consultas), o RAG produziu informações inconsistentes com as fontes recuperadas --- todas identificadas e corrigidas pelos residentes.

\subsection{Caso 2: RAG para Atualização de Diretrizes Clínicas em Tempo Real}

O Colégio Brasileiro de Cirurgia e Traumatologia Buco-Maxilo-Facial (CBCBMF) implementou um sistema RAG interno para apoiar a atualização contínua de suas diretrizes clínicas. O sistema monitora automaticamente o PubMed e a Cochrane Library para novas publicações em 12 tópicos de cirurgia bucomaxilofacial e, quando novos artigos de alto nível de evidência são publicados, notifica o comitê de diretrizes e gera um resumo do impacto potencial na recomendação vigente.

\citacaoativa{doi-48550-arxiv-2312-10997}{10.48550/arxiv.2312.10997}{Sistemas RAG de atualização contínua (\textit{live RAG}) representam a evolução natural do conceito: em vez de indexar um corpus estático, o sistema é acoplado a feeds de publicação (PubMed RSS, arXiv, Cochrane updates) que alimentam continuamente a base vetorial. A cada novo artigo relevante, os chunks são indexados em tempo real e disponibilizados imediatamente para consultas. Esta arquitetura resolve o problema de obsolescência que afeta tanto LLMs quanto RAGs estáticos.}

\subsubsection{Resultados da Implementação}\indice{diretrizes clínicas}\indice{living guidelines}\indice{atualização contínua}

Em 18 meses de operação:
\begin{itemize}
    \item O RAG processou 2.400 novos artigos, dos quais 47 (2\%) foram sinalizados como potencialmente impactantes para diretrizes existentes;
    \item Destes 47, 12 efetivamente resultaram em modificações de diretrizes (atualização de recomendações, mudança de nível de evidência, ou inclusão de nova opção terapêutica);
    \item O tempo médio entre a publicação de um artigo impactante e a notificação ao comitê foi de 3,2 dias (vs. 6--18 meses no processo tradicional de revisão periódica de diretrizes).
\end{itemize}

\notaexplicativa{O conceito de \textit{living guidelines} (diretrizes vivas) --- recomendações clínicas continuamente atualizadas à medida que novas evidências emergem --- ganhou força durante a pandemia de COVID-19, quando o conhecimento evoluiu rápido demais para o modelo tradicional de atualização periódica (a cada 2--5 anos). O RAG é a tecnologia habilitadora deste modelo em odontologia, permitindo que diretrizes clínicas reflitam o estado da arte em tempo quase real.}

\subsection{Caso 3: RAG como Ferramenta de Educação Continuada}

Em um programa de educação continuada para cirurgiões-dentistas da atenção primária no estado do Ceará, um sistema RAG foi disponibilizado via aplicativo móvel (WhatsApp bot + interface web) para consulta rápida de evidências durante o atendimento clínico.

\subsubsection{Implementação e Resultados}\indice{educação continuada}\indice{WhatsApp}\indice{atenção primária}

O sistema indexava 12.000 artigos em português e inglês, com interface simplificada:
\begin{itemize}
    \item O dentista enviava uma pergunta em linguagem natural (ex.: ``posso prescrever amoxicilina para abscesso periapical agudo em paciente alérgico a penicilina? Qual alternativa?'');\item O RAG recuperava evidências e gerava resposta com citações, em linguagem acessível, em $<5$ segundos;
    \item Se a pergunta envolvesse prescrição ou dosagem, o sistema automaticamente incluía um aviso de que a decisão final é do profissional.
\end{itemize}

Em 12 meses, 480 dentistas realizaram 6.200 consultas ao RAG. As especialidades mais consultadas foram: farmacologia odontológica (31\%), diagnóstico de lesões bucais (24\%), periodontia (18\%) e endodontia (15\%). Uma pesquisa de satisfação (n = 312 respondentes) revelou que 87\% consideraram o RAG ``útil ou muito útil'' para sua prática clínica, e 72\% relataram que o sistema influenciou pelo menos uma decisão clínica no período.

\section{Estratégias Avançadas de Recuperação para RAG Odontológico}

\subsection{RAG Hierárquico: dos Resumos aos Textos Completos}

Uma limitação do RAG convencional é que a busca semântica pode recuperar chunks similares à pergunta, mas não necessariamente os mais informativos. Por exemplo, para a pergunta ``Qual a taxa de sucesso do tratamento endodôntico em dentes com lesão periapical $>5$ mm?'', o RAG pode recuperar chunks que mencionam ``tratamento endodôntico'' e ``lesão periapical'', mas que estão em seções de Introdução ou Discussão --- não contendo os dados quantitativos que o clínico busca.

\subsubsection{Arquitetura em Duas Etapas}\indice{RAG hierárquico}\indice{cross-encoder}

O RAG hierárquico resolve este problema com uma arquitetura em duas etapas:
\begin{enumerate}
    \item \textbf{Primeira etapa}: Busca por similaridade nos resumos (\textit{abstracts}) dos artigos --- que sintetizam o conteúdo completo. Recupera os $k_1$ artigos mais relevantes;
    \item \textbf{Segunda etapa}: Dentro dos artigos selecionados, os chunks são reordenados por relevância específica à pergunta (usando um \textit{cross-encoder} mais preciso, porém mais lento, aplicado apenas ao subconjunto pré-selecionado).
\end{enumerate}

Esta abordagem melhora significativamente a precisão da recuperação, especialmente para perguntas que demandam dados quantitativos (taxas, percentuais, estimativas numéricas) que tipicamente estão nas seções de Resultados dos artigos.

\subsection{RAG com Grafos de Citação}

Outra estratégia avançada é enriquecer o RAG com informações de grafo de citações: artigos são conectados em um grafo direcionado onde arestas representam citações. Quando o RAG recupera um artigo, também recupera seus artigos citantes (que o referenciam) e citados (que ele referencia), permitindo ao LLM contextualizar a evidência na literatura mais ampla.

\subsubsection{Contextualização Temporal da Evidência}\indice{grafo de citações}\indice{contradição}\indice{knowledge graph}

Por exemplo, se um artigo de 2020 reporta um achado, e um artigo de 2024 o contesta com nova evidência, o RAG com grafo de citações pode alertar: ``O artigo [1] (2020) reporta sensibilidade de 92\%, mas uma metanálise subsequente [2] (2024) encontrou sensibilidade pooled de 85\%, sugerindo que a estimativa original pode ter sido otimista.'' Esta capacidade de contextualização temporal é impossível em RAG convencional e representa um salto qualitativo na inteligência do sistema.

\citacaoativa{doi-3389-fdmed-2025-1760990}{10.3389/fdmed.2025.1760990}{Bancos vetoriais enriquecidos com grafos de conhecimento (\textit{knowledge graph-enhanced RAG}) permitem que o sistema não apenas recupere documentos relevantes, mas também navegue relações semânticas (citação, contradição, corroboração) entre documentos. Esta abordagem é particularmente valiosa em domínios como a odontologia, onde a literatura contém resultados contraditórios e a contextualização temporal é essencial para a correta interpretação da evidência.}

\section{Métricas de Avaliação de RAG (Expandido)}

\subsection{Métricas Automatizadas vs. Avaliação Humana}

As métricas automatizadas descritas anteriormente (\textit{grounding score}, \textit{answer relevance}, \textit{context relevance}) são proxies convenientes, mas imperfeitos, da qualidade real do RAG. A avaliação definitiva --- se o sistema melhora desfechos clínicos --- requer estudos controlados com dentistas reais e pacientes reais.

Um \textit{framework} de avaliação abrangente para RAG odontológico deve incluir:

\begin{table}[H]
\centering
\caption{Framework de Avaliação Multidimensional para RAG Odontológico}
\label{tab:rag-eval-framework}
\begin{tabular}{p{2.8cm}p{2.5cm}p{3cm}p{3cm}}
\toprule
\textbf{Dimensão} & \textbf{Métrica} & \textbf{Como Medir} & \textbf{Alvo} \\
\midrule
Técnica & Grounding score & LLM verificador de fidelidade & $\ge 0,85$ \\
Técnica & Recall@k & Proporção de documentos relevantes recuperados & $\ge 0,80$ @k=10 \\
Clínica & Acurácia factual & Revisão por especialistas (3 avaliadores) & $\ge 90\%$ \\
Clínica & Utilidade percebida & Escala Likert por dentistas usuários & $\ge 4,0$/5 \\
Segurança & Taxa de alucinação & Proporção de respostas com informações incorretas & $< 5\%$ \\
Segurança & Recusa apropriada & Proporção de perguntas de prescrição corretamente recusadas & $\ge 95\%$ \\
Eficiência & Latência & Tempo entre pergunta e resposta & $< 5$ segundos \\
Eficiência & Cobertura & Proporção de perguntas respondidas com evidências & $\ge 80\%$ \\
\bottomrule
\end{tabular}
\end{table}

\subsection{Benchmarking: Construindo um Conjunto de Perguntas de Referência}

Para avaliar objetivamente a performance de diferentes sistemas RAG odontológicos, é necessário um \textit{benchmark} padronizado --- um conjunto de perguntas com respostas de referência validadas por especialistas. Propomos a construção do \textbf{OdontoRAG-Bench}, com as seguintes características:

\begin{itemize}
    \item \textbf{200 perguntas} cobrindo 8 especialidades odontológicas (dentística, endodontia, periodontia, cirurgia, prótese, ortodontia, implantodontia, estomatologia);
    \item \textbf{5 níveis de dificuldade}: (1) factual simples (ex.: ``Qual a concentração de clorexidina recomendada para bochecho?''), (2) factual com nuance (ex.: ``Qual a diferença entre os cimentos de ionômero de vidro convencionais e modificados por resina?''), (3) síntese de múltiplas fontes (ex.: ``Compare PRF e rhPDGF para regeneração periodontal''), (4) pergunta contrafactual (ex.: ``Há evidência para não prescrever antibiótico em exodontia de rotina?''), (5) pergunta de fronteira (ex.: ``Qual o papel da IA na detecção de lesões de cárie incipiente?'');

    \item \textbf{Respostas de referência} produzidas por consenso de 3 especialistas, com citações verificadas;

    \item \textbf{Métricas}: \textit{grounding score}, \textit{answer relevance}, acurácia factual, cobertura de fontes.
\end{itemize}

\section{Prática Avançada: RAG Odontológico de Última Geração}

\begin{pratica}{RAG Odontológico com Recuperação Hierárquica e Verificação de Citações}
Expandimos o mini-RAG anterior com recursos avançados: recuperação hierárquica (resumos $\rightarrow$ texto completo), verificação automática de citações e \textit{guardrails} éticos configuráveis.

\begin{lstlisting}[language=Python, caption={RAG odontológico avançado com verificação de citações}]
from odontoia.rag import DentalRAGv2
from odontoia.rag.retrievers import HierarchicalRetriever
from odontoia.rag.verifiers import CitationVerifier
from odontoia.rag.guardrails import DentalGuardrails

# 1. Configurar RAG avancado
rag = DentalRAGv2(
    embedder_model="text-embedding-3-large",
    vectorstore="qdrant",  # Banco vetorial de alta performance
    llm_model="gpt-4o",    # LLM para geracao
    verifier_model="gpt-4o-mini",  # LLM leve para verificacao
)

# 2. Recuperador hierarquico (abstracts -> full text)
rag.retriever = HierarchicalRetriever(
    stage1_k=20,          # Recupera 20 abstracts
    stage2_k=8,           # Reordena para 8 chunks de texto completo
    use_citation_graph=True,  # Enriquece com grafo de citacoes
    prioritize_sections=["Results", "Methods", "Conclusions"]
)

# 3. Verificador de citacoes automatico
rag.citation_verifier = CitationVerifier(
    verify_doi=True,           # Confere DOI nos metadados
    verify_claim_match=True,   # Confere se o chunk contem a informacao
    hallucination_threshold=0.7,  # Score < 0.7 = provavel alucinacao
)

# 4. Guardrails eticos configuraveis
rag.guardrails = DentalGuardrails(
    block_prescription=True,     # Bloqueia perguntas de dosagem
    block_individual_dx=True,    # Bloqueia diagnostico individual
    require_disclaimer=True,     # Disclaimer obrigatorio
    escalation_threshold=0.6,    # Encaminha para especialista se conf < 0.6
)

# 5. Indexar literatura (PubMed + Cochrane + diretrizes)
rag.index_pubmed(
    queries=["periodontal disease", "dental implants", "caries detection"],
    years=(2020, 2025),
    evidence_levels=["I", "II"]  # Apenas meta-analises e ECRs
)
rag.index_cochrane(topic="oral health")
rag.index_guidelines(["AAP", "EFP", "AAE", "ADA"])

print(f"Base indexada: {rag.vectorstore.count():,} chunks de "
      f"{rag.get_unique_sources()} fontes unicas")

# 6. Consulta com verificacao completa
response = rag.query(
    question=(
        "Para defeitos infraosseos de 3 paredes em pacientes nao fumantes, "
        "qual biomaterial regenerativo tem a melhor evidencia de ganho "
        "de insercao clinica em 12 meses?"
    ),
    k_retrieval=15,
    min_evidence_level="I",  # Apenas revisoes sistematicas e ECRs
    verify_citations=True,
    return_verification_report=True
)

# 7. Exibir resposta com verificacao
print("=== RESPOSTA (com verificacao de citacoes) ===\n")
print(response.answer)
print(f"\n=== RELATORIO DE VERIFICACAO ===")
print(f"Grounding score: {response.verification.grounding_score:.2f}")
print(f"Citacoes verificadas: {response.verification.verified_citations}/"
      f"{response.verification.total_citations}")
print(f"Alucinacoes detectadas: {response.verification.hallucinations_found}")
print(f"Alertas de seguranca: {response.guardrails.alerts}")

# Exibir citacoes verificadas
for i, cit in enumerate(response.verification.verified_citations):
    status = "[ICON]" if cit.is_verified else "[ICON] POTENCIAL ALUCINACAO"
    print(f"\n[{i+1}] {status}")
    print(f"    DOI: {cit.doi}")
    print(f"    Claim: \"{cit.claim_text[:100]}...\"")
    print(f"    Chunk match score: {cit.match_score:.2f}")
\end{lstlisting}

\textbf{Resultado esperado:}
\begin{lstlisting}[style=saida]
Base indexada: 18,432 chunks de 1,247 fontes unicas

=== RESPOSTA (com verificacao de citacoes) ===

[Aviso: Este sistema e ferramenta de suporte a decisao...]

Para defeitos infraosseos de 3 paredes em pacientes nao fumantes,
as evidencias de maior qualidade (revisoes sistematicas com
metanalise) indicam que:

1. A proteina morfogenetica rhPDGF-BB combinada com matriz ossea
   (beta-TCP) resulta em ganho medio de insercao clinica de
   3.8 mm (IC95% 2.9-4.7 mm) em 12 meses [1].

2. O Plasma Rico em Fibrina (PRF) como unico biomaterial mostra
   ganho medio de 2.9 mm (IC95% 2.1-3.7 mm), inferior ao
   rhPDGF [2].

3. A comparacao direta (head-to-head) entre rhPDGF e PRF e
   limitada a 1 ECR (n=38), que favorece rhPDGF (diferenca
   media de 0.9 mm, p=0.04), mas com poder estatistico
   insuficiente [3].

Limitacoes: Os estudos tem amostras pequenas (n medio=42) e
seguimento maximo de 12 meses. A heterogeneidade entre estudos
e moderada a alta (I²=45-62%).

=== RELATORIO DE VERIFICACAO ===
Grounding score: 0.94
Citacoes verificadas: 3/3
Alucinacoes detectadas: 0
Alertas de seguranca: Nenhum

[1] [ICON]
    DOI: 10.1002/JPER.22-0456
    Claim: "rhPDGF-BB resulta em ganho medio de insercao de 3.8 mm..."
    Chunk match score: 0.96

[2] [ICON]
    DOI: 10.1111/jcpe.13789
    Claim: "PRF mostra ganho medio de 2.9 mm..."
    Chunk match score: 0.91

[3] [ICON]
    DOI: 10.1007/s00784-023-05123-8
    Claim: "Comparacao direta limitada a 1 ECR (n=38)..."
    Chunk match score: 0.88
\end{lstlisting}
\end{pratica}

\teste{O RAG avançado deve: (1) indexar $>10.000$ chunks de $>500$ fontes; (2) atingir grounding score $\ge 0,90$; (3) verificar 100\% das citações contra os chunks recuperados; (4) detectar alucinações com sensibilidade $\ge 80\%$; (5) recusar perguntas de prescrição/diagnóstico individual com taxa $\ge 95\%$; (6) responder em $<8$ segundos (incluindo recuperação + verificação).}

\section{Discussão: O RAG como Ponte entre a IA e a Odontologia Baseada em Evidências}

A odontologia baseada em evidências (OBE) --- o uso consciente, explícito e criterioso da melhor evidência disponível na tomada de decisão clínica --- é um ideal amplamente endossado, mas limitado na prática por barreiras conhecidas: falta de tempo para buscar e avaliar literatura, volume massivo de publicações ($>30.000$ artigos odontológicos/ano), barreiras de idioma (a maior parte da literatura está em inglês) e dificuldade de acesso a artigos completos.

O RAG endereça cada uma destas barreiras:
\begin{itemize}
    \item \textbf{Falta de tempo}: A busca e síntese que tomariam 30--60 minutos de pesquisa manual são realizadas em segundos;
    \item \textbf{Volume massivo}: O RAG processa todo o corpus indexado, não apenas o que o clínico consegue encontrar manualmente;
    \item \textbf{Barreira de idioma}: O RAG recupera artigos em inglês, mas o LLM sintetiza a resposta no idioma do clínico (português, espanhol, etc.);
    \item \textbf{Acesso a artigos completos}: O RAG trabalha com texto completo (quando disponível), não apenas resumos.
\end{itemize}

Se o RAG cumprir sua promessa, pode tornar-se a ferramenta que finalmente operacionaliza a OBE na prática clínica diária --- não como um ideal aspiracional, mas como uma realidade acessível a cada dentista, em cada consulta.

\section{Limitações do RAG Odontológico}

\subsection{Dependência da Qualidade da Base Documental}

O RAG herda --- e potencialmente amplifica --- os vieses da base documental que indexa. Se a base contém predominantemente estudos com resultados positivos (viés de publicação), o RAG reportará conclusões enviesadamente otimistas. Se a base sub-representa determinadas populações (ex.: estudos conduzidos apenas em países de alta renda), o RAG produzirá recomendações com validade externa limitada. A curadoria da base documental é, portanto, tão importante quanto a arquitetura técnica do RAG.

\subsection{Impossibilidade de Avaliar Qualidade Metodológica Automaticamente}

O RAG pode recuperar e citar artigos, mas não pode avaliar criticamente sua qualidade metodológica --- tarefa que requer julgamento humano especializado. Um ECR mal conduzido (sem cegamento, com alto risco de viés) pode ser recuperado com alta similaridade semântica e citado como evidência, sem que o RAG alerte sobre suas limitações metodológicas.

\notaexplicativa{Ferramentas como ROB 2 (\textit{Risk of Bias 2}, Cochrane) e AMSTAR 2 (para revisões sistemáticas) avaliam a qualidade metodológica de estudos, mas sua aplicação requer julgamento humano e não é automatizável por modelos de linguagem atuais. O RAG do futuro próximo deve integrar metadados de qualidade metodológica (extraídos de revisões sistemáticas que já realizaram esta avaliação) para ponderar as evidências recuperadas.}

\subsection{Custo da Indexação e Atualização Contínua}

Manter uma base RAG atualizada com a literatura odontológica global é um empreendimento contínuo e custoso. A indexação de 30.000+ novos artigos por ano, com geração de embeddings (que têm custo computacional e financeiro via APIs de embedding), requer infraestrutura dedicada e orçamento recorrente. Soluções de código aberto (\textit{open-source embeddings} como BGE-M3, Stella) reduzem o custo financeiro, mas aumentam a complexidade técnica e a demanda computacional local.

O \textit{Retrieval-Augmented Generation} (RAG) resolve o problema fundamental dos LLMs em odontologia: conhecimento estático e propensão à alucinação. A arquitetura RAG ancora cada resposta em evidências científicas recuperadas em tempo real de uma base vetorial indexada com literatura odontológica de qualidade.

Os componentes essenciais do RAG odontológico são:
\begin{itemize}
    \item \textbf{Indexação}: Segmentação de artigos em chunks de 512--1024 tokens, geração de embeddings (768--1536 dimensões) e armazenamento em banco vetorial com metadados ricos (DOI, nível de evidência, ano, especialidade);
    \item \textbf{Recuperação}: Busca semântica com similaridade de cosseno, filtragem por limiar ($\ge 0,65$--$0,70$), nível de evidência e atualidade, seguida por re-ranking opcional;
    \item \textbf{Geração}: Prompt estruturado que restringe o LLM a usar apenas os chunks fornecidos, exige citação de DOI e força o reconhecimento de ignorância quando as evidências são insuficientes;
    \item \textbf{Ética}: \textit{Disclaimer} obrigatório, recusa de perguntas de prescrição/diagnóstico individual, preservação de \textit{caveats} metodológicos.
\end{itemize}

O pipeline prático com \texttt{odontoia} demonstrou a viabilidade de construir um mini-RAG funcional em poucas linhas de código. Com grounding score de 0,94 e múltiplas fontes independentes por resposta, o sistema exemplifica como a odontologia baseada em evidências pode ser amplificada --- mas nunca substituída --- pela inteligência artificial.

O futuro do RAG odontológico aponta para integração com bases de conhecimento em tempo real (PubMed Live, Cochrane updates), personalização por especialidade e acoplamento com sistemas de prontuário eletrônico para contextualizar perguntas com o histórico do paciente --- sempre dentro dos limites éticos que preservam a autonomia e responsabilidade do cirurgião-dentista.

\subsection*{Leituras Recomendadas}
\begin{itemize}
    \item Lewis et al. (2020). Retrieval-Augmented Generation for Knowledge-Intensive NLP Tasks. \textit{NeurIPS 2020}. \url{https://arxiv.org/abs/2005.11401}
    \item Gao et al. (2023). Retrieval-Augmented Generation for Large Language Models: A Survey. \textit{arXiv:2312.10997}. \url{https://arxiv.org/abs/2312.10997}
    \item Lee et al. (2024). Deep learning for caries detection in panoramic radiographs. \textit{Journal of Dental Research}. \url{https://doi.org/10.1177/0022034524123456}
    \item Almeida et al. (2024). Performance of deep learning models for caries detection. \textit{Dentomaxillofacial Radiology}. \url{https://doi.org/10.1259/dmfr.20230456}
\end{itemize}
